\documentclass[11pt]{article}

% basic packages
\usepackage[margin=1in]{geometry}
\usepackage[pdftex]{graphicx}
\usepackage{amsmath,amssymb,amsthm}
\usepackage{custom}
\usepackage{lipsum}
\usepackage{booktabs}
\usepackage{enumitem}

% page formatting
\usepackage{fancyhdr}
\pagestyle{fancy}

\renewcommand{\sectionmark}[1]{\markright{\textsf{\arabic{section}. #1}}}
\renewcommand{\subsectionmark}[1]{}
\lhead{\textbf{\thepage} \ \ \nouppercase{\rightmark}}
\chead{}
\rhead{}
\lfoot{}
\cfoot{}
\rfoot{}
\setlength{\headheight}{14pt}

\linespread{1.03} % give a little extra room
\setlength{\parindent}{0.2in} % reduce paragraph indent a bit
\setcounter{secnumdepth}{2} % no numbered subsubsections
\setcounter{tocdepth}{2} % no subsubsections in ToC

\begin{document}

% make title page
\thispagestyle{empty}
\bigskip \
\vspace{0.1cm}

\begin{center}
{\fontsize{22}{22} \selectfont Lecture Notes on}
\vskip 16pt
{\fontsize{36}{36} \selectfont \bf \sffamily Quantum Mechanics}
\vskip 24pt
{\fontsize{18}{18} \selectfont \rmfamily Owen Lin} 
\vskip 6pt
{\fontsize{14}{14} \selectfont \ttfamily owen2006lin@gmail.com} 
\vskip 24pt
\end{center}

{\parindent0pt \baselineskip=15.5pt
These notes cover quantum mechanics at the undergraduate level. The sources used are:
\begin{itemize}
    \item Griffiths: \textit{Introduction to Quantum Mechanics}
\end{itemize}


Potential sources to consult in the future:

\begin{itemize}
    \item Sakurai: \textit{Modern Quantum Mechanics}
\end{itemize}

}
 


% make table of contents
\newpage
\microtoc
\newpage

% main content
\section{The Schr\"odinger Equation}

\subsection{Wave Functions}
Classically, we work with masses and attempt to solve for their equations of motions \( x(t) \) using corresponding
Newton's laws or energy equations. The quantum mechanics approach is different; the idea is to instead solve for a
particle's wave function \( \Psi(x,t) \) and from there we can do anything. The goal then is to solve the Schr\"odinger
equation

\[ 
    i\hbar \frac{\partial \Psi}{\partial t} = -\frac{\hbar}{2m}\frac{\partial^{2} \Psi}{\partial x^{2}} + V\Psi
\]

\noindent After solving the Schr\"odinger equation for some intial wave function \( \Psi(x, 0) \) the equation actually gives us 
all future values of the wave function \( \Psi(x, t) \).

\noindent The statistical interpretation of the wave function goes as follows:
\begin{itemize}
    \item Born's statistical interpretation of the wave function states that 
    
    \[ 
        \int^{a}_{b} |\Psi(x,t)|^{2}dx  
    \]

    is the probability of finding the particle at \( a \leq x \leq b \) at time t. Note that while \( \Psi \) is a complex valued
    function, \( |\Psi|^{2} = \Psi \Psi^{*}\) where \( \Psi^{*} \) is the complex conjugate of \( \Psi \) is always real and nonnegative.

    \item An obvious consequence is this is that even if you know everything about a particle (ie; you know its wavefunction), you still
    cannot predict with certainty the outcome of an experiment to measure its exact position. 

    \item Before measurement, a particle's position is not well defined. Only after measurement can we say the particle has a "position," and
    by doing so we collapse the particle's wavefunction to a dirac delta centered around the measured position. This is known as the\textbf{
    Copenhagen Interpretation}

    \item The integral over all possible positions must be equal to one, since the particle must be located somewhere.
    \[ 
        \int^{\infty}_{-\infty} |\Psi(x,t)|^{2}dx = 1
    \]

    If we have a wavefunction, we need to make sure its normalized for it to be physical. We can normalize it to obey the integral by multiplying
    it by a scale factor. A glance at the Schr\"odinger equation one can see that if \( \Psi(x,t) \) is a solution, then \( A\Psi(x,t) \) is also
    a solution for nonzero constant A.

    \item One can also show that the normalization integral does not change with time, so there is no need to constantly renormalized solutions as 
    the particle evolves with time.

    \item If we interpret \( |\Psi(x,t)|^{2} \) to be the probability, then we can say the expectation value of position is 
    \[ 
        \langle x \rangle = \int^{\infty}_{-\infty} x|\Psi(x,t)|^{2}dx
    \]
    Note that the interpretation of this quantity is the expected value of the position from measuring particles with wave function \( \Psi \). Simply
    remeasuring the same particle over again will not work, since the first measurement will collapse the particle into some other wave function with
    exact position

    \item By differentiating the expected position with respect to time, one can obtain the expectation value of the velocity. It's customary for us to
    work using momentum so we have
    
    \[ 
        \langle p \rangle = -i\hbar \int \left(\Psi^{*}\frac{\partial \Psi}{\partial x}\right)dx 
    \]

    \item A more illustrative way to write the expectation values for position and momentum are:
    \[ 
        \langle x \rangle = \int \Psi^{*} \cdot x \cdot \Psi dx
    \]

    \[ 
        \langle p \rangle = -i\hbar \int \Psi^{*} \cdot \left(\frac{\partial}{\partial x}\right) \cdot \Psi dx 
    \]
    We say the the \( x \) operator "represents" position and the \( -i \hbar\frac{\partial}{\partial x} \)  operator "represents" momentum. 

    \item All classical quantities can be represented as a function of position and angular momentum. Then the representation of a classical quantity \( Q(x,p) \)
    can be written as 
    \[ 
        \langle Q(x,p) \rangle = \int \Psi^{*} Q\left(x, -i\hbar\frac{\partial}{\partial x}\right) \Psi dx. 
    \]

    \item \textbf{Ehrenfest's Theorem} states that expectation values obey their classical laws. For example, the time derivative of momentum is
    
    \[ 
        \frac{d\langle p \rangle}{dt} = \langle - \frac{\partial V}{\partial x}\rangle 
    \]

\end{itemize}

\subsection{Time Independent Schr\"odinger Equation}
Let us assume that the potential \(V(x)\) is independent of time. Then using separation of variables, we can substitute \(\Psi(x,t) = \psi(x)\phi(t)\)
into the Schr\"odinger equation giving us 
\(i\hbar\frac{1}{\phi}\frac{d\phi}{dt} = -\frac{\hbar^2}{2m} \frac{1}{\psi}\frac{d^2\psi}{dx^2} + V\). Then for the equality to hold, both sides must be equal to 
some constant value \(E\) leaving us with two ordinary differential equations:
\[
    \frac{d\phi}{dt} = -\frac{iE}{\hbar}\phi
\]

\[
    -\frac{\hbar^2}{2m}\frac{1}{\psi}\frac{d^2\psi}{dx^2} + V = E
\]
The second equation is the \textbf{Time Independent Schr\"odinger Equation}. The first equation defines the characteristic time dependence and can
be solved trivially \(\phi(t) = e^{-iEt/\hbar}\). Separable solutions obey several important properties that we're interested in:

\begin{itemize}
    \item These solutions are \textbf{stationary states}. What this means is while the wave function is obviously time dependent, 
    \(\Psi(x,t) = \psi(x)e^{-iEt/\hbar}\), the probability density is not:
    \[
        |\Psi(x,t)|^2 = \Psi \Psi^* = \psi^* e^{+iEt/\hbar}\psi e^{-iEt/\hbar} = |\psi(x)|^2
    \]
    We can further show that the expectation value of any classical quantity is also not time dependent in a stationary state:
    \[ 
        \langle Q(x,p) \rangle = \int \psi^{*} Q\left(x, -i\hbar\frac{\partial}{\partial x}\right) \psi dx. 
    \]


    \item The constant \(E\) represents the Hamiltonian of the system. Classically, the Hamiltonian is \(H(x,p) = \frac{p^2}{2m} + V(x)\) so the 
    corresponding Hamiltonian operator is \(\hat{H} = -\frac{\hbar^2}{2m}\frac{\partial^2}{\partial x^2} + V(x)\). So the time independent Schr\"odinger
    equation can be rewritten as 
    \[
        \hat{H}\psi = E\psi
    \]

    Then the expectation values of the Hamiltonian is 

    \[
        \langle H \rangle = \int \psi^{*} \hat{H} \psi dx = \int \psi^{*} E \psi dx = E
    \]

    \[
        \langle H^2 \rangle = \int \psi^{*} \hat{H}^2 \psi dx = \int \psi^{*} E (\hat{H}\psi) dx = \int \psi^{*} E^2 \psi dx= E^2
    \]
    This means the standard deviation is 0. So any ensemble of stationary states with the same wavefunction have their Hamiltonian measured to be the exact
    same E. 

    \item Every solution to the Schr\"odinger equation can be written in the form of a sum of stationary solutions. 
    \[
        \Psi(x,t) = \sum_{n=1}^{\infty}c_n \psi_n(x)e^{-E_n t / \hbar}     
    \]
    
    So we can basically decompose any wave function into the sum of wavefunctions associated with specific energy levels. The quantity \(|c_n|^2\) represents 
    the probability that a measurement of the energy will return value \(E_n\), so we also must demand 
    \[
        \sum_{n=1}^{\infty} |c_n|^2 = 1 
    \]
    \[
        \langle H \rangle = \sum_{n=1}^{\infty} |c_n|^2E_n 
    \]
    

\end{itemize}

\subsection{Infinite Potential Well}
\noindent Consider the potential function given by 
    \[
    V(x) = \begin{cases}
        0 & 0 \leq x \leq a \\
        \infty & otherwise
    \end{cases}        
    \]

\noindent Outside the potential well, \(\psi(x)\) must be 0 because \(V(x)\) goes to infinity. So we have \(\psi(x) = 0\) for \(x < 0\) amd \(x > a\). Inside
the well, the Schr\"odinger equation gives us \(-\frac{\hbar}{2m}\frac{d^2\psi}{dx^2} = E\psi\). The solution is just the simple harmonic oscillator with 
\(k = \sqrt{2mE}/\hbar\). So the general solution is \(\psi(x) = Asin(kx) + Bcos(kx)\). Using the boundary condition that the wavefunction must go to 0 at
\(x = 0\) and \(x = a\), we obtain
\[
    \psi_n(x) = \sqrt{\frac{2}{a}}\sin{\left(\frac{n\pi}{a}x\right)}
\]
with corresponding energy 
\[
    E_n = \frac{n^2\pi^2\hbar^2}{2ma^2}
\]

\noindent Some properties of the \(\psi_n\) functions:
\begin{itemize}
    \item \(\psi_n\) alternates parity around the center of the well. \(\psi_1\), the ground state, is even, \(\psi_2\), the first excited state is odd, \(\psi_3\)
    is even and so on.
    \item Each \(\psi_n\) are orthonormal. Namely:
     \[
        \int \psi_m(x)* \psi_n(x) dx = \delta_{mn}    
    \]
    \item They form a complete basis for any function \(f(x)\):
    \[
        f(x) = \sum_{n=1}^{\infty} c_n \psi_n(x) = \sqrt{\frac{2}{a}}\sum_{n= 1}^{\infty} c_n \sin{\left(\frac{n\pi}{a}x\right)}     
    \]
    This is simply a fourier series for \(f(x)\), and this can represent any function \(f(x)\) by \textbf{Dirichlet's Theorem}.
    To actually determine the coefficients, multiply both sides by \(\psi_m^*\) and then integrate
    \[
        \int \psi_m(x)^* f(x) dx = \sum_{n=1}^{\infty} c_n\int \psi_m(x)^* \psi_n(x)dx = \sum_{n=1}^{\infty}c_n \delta_{nm} = c_m 
    \]
    \[
        c_n = \int \psi_n(x)^* f(x)dx
    \]
    
    
    
\end{itemize}





\subsection{The Quantum Harmonic Oscillator}

\section{Dirac Notation}

\section{The Hydrogen Atom}

\section{Symmetries and Conservation Laws}

\end{document}